\documentclass[11pt]{amsart}
\usepackage{geometry}                % See geometry.pdf to learn the layout options. There are lots.
\geometry{letterpaper}                   % ... or a4paper or a5paper or ... 
%\geometry{landscape}                % Activate for for rotated page geometry
%\usepackage[parfill]{parskip}    % Activate to begin paragraphs with an empty line rather than an indent
\usepackage{graphicx}
\usepackage{amssymb}
\usepackage{epstopdf}
\DeclareGraphicsRule{.tif}{png}{.png}{`convert #1 `dirname #1`/`basename #1 .tif`.png}

\title{PS 2.2}
\author{Viridiana Josselyne}
%\date{}                                           % Activate to display a given date or no date

\begin{document}
\maketitle
%\section{}
%\subsection{}
\textbf{PS 2.2} \newline

La cotangente de un ángulo se define como el cociente entre el coseno y el seno de dicho ángulo.

\begin{center}

	cot g(a) = $ \frac{cos(a)}{sen(a)} $ \newline
	\newline
	\textbf{Fórmula 6.21} \newline

\end{center}  

Construya un diagrama de flujo que le permita calcular la cotangente de un ángulo, considerando que se conoce el valor del seno y del coseno del mismo. Recuerde que el seno debe ser diferente de 0. \newline

Datos: SENO, COSENO (variables de tipo real).

\end{document}  